\chapter{\ifproject%
\ifenglish Experimentation and Results\else การทดลองและผลลัพธ์\fi
\else%
\ifenglish System Evaluation\else การประเมินระบบ\fi
\fi}

ในบทนี้นำเสนอการทดลองเบื้องต้นที่ดำเนินการแล้ว ได้แก่
(1) การทดลองส่ง Raw Logs ให้ LLM เชิงพาณิชย์วิเคราะห์โดยตรง เพื่อยืนยันปัญหาของโครงงาน
(2) การทดลองเพื่อตอบคำถามว่า LLM ขนาดเล็กยังจำเป็นต้องใช้ RAG หรือไม่
และ (3) แผนการทดลองและการประเมินผลในขั้นต่อไป

\section{การทดลองเบื้องต้น: ส่ง Raw Logs ให้ LLM วิเคราะห์โดยตรง}

\subsection{วิธีการทดลอง}
ดึง Windows Sysmon Logs ของชุดข้อมูล T1087.001 จาก Splunk ผ่าน API ได้เป็นไฟล์ JSON
แล้วอัปโหลดไฟล์ให้ LLM เชิงพาณิชย์ 3 โมเดลวิเคราะห์โดยตรง
ด้วยคำถามว่า ``นี่คือ Attack Technique อะไร เพราะอะไรถึงคิดแบบนั้น''
คำตอบที่คาดหวังคือ Account Discovery: Local Account (T1087.001)

\subsection{ผลการทดลอง}
\begin{table}[ht]
  \centering
  \caption{ผลการส่ง Raw Logs ของ T1087.001 ให้ LLM เชิงพาณิชย์วิเคราะห์โดยตรง}
  \label{tab:raw-llm}
  \begin{tabular}{|L{0.25\textwidth}|L{0.45\textwidth}|L{0.17\textwidth}|}
    \hline
    \textbf{โมเดล} & \textbf{คำตอบ} & \textbf{ผล} \\ \hline
    GPT-5.6 Luna      & T1003 OS Credential Dumping & ผิด \\ \hline
    Gemini 3.1 Pro    & ไม่พบพฤติกรรมการโจมตีที่ชัดเจน (Logs เป็นเพียงการทำงานของ Splunk Agent) & ผิด \\ \hline
    Claude Sonnet 5 (High) & T1087 Account Discovery & ถูกระดับ Technique \\ \hline
  \end{tabular}
\end{table}

จากตารางที่~\ref{tab:raw-llm} พบว่าแม้เป็นโมเดลขนาดใหญ่ แต่เมื่อได้รับ Raw Logs ที่มี Noise จำนวนมาก
บางโมเดลถูกดึงความสนใจไปที่ Events ที่ไม่เกี่ยวข้อง (เช่น การเข้าถึง Process ของ Splunk Agent)
และไม่มีโมเดลใดระบุถึงระดับ Sub-technique ได้ถูกต้อง
ผลนี้สนับสนุนสมมติฐานว่าการสร้าง Investigation Context ที่คัดกรองแล้วมีความจำเป็น

\section{การทดลอง: ความจำเป็นของ RAG สำหรับ LLM ขนาดเล็ก}

\subsection{วัตถุประสงค์และการตั้งค่า}
เพื่อทดสอบว่า เมื่อผ่าน Event Selection แล้ว LLM ขนาดเล็ก (Qwen3.5:9B) ยังต้องใช้ RAG หรือไม่
และ Query แบบใดเหมาะกับการค้นหาในฐานความรู้
Input ของการทดลองคือผลลัพธ์จาก Event Selection แบบ Top-$k$ ของชุดข้อมูล T1087.001
ประกอบด้วย Alert และ 16 Events (EID 1 จำนวน 3 Events ได้แก่ \texttt{cmd.exe} $\rightarrow$ \texttt{net users}
$\rightarrow$ \texttt{net1 users} และที่เหลือเป็น EID 10) ในรูปแบบ JSON
เงื่อนไขการทดลองแสดงในตารางที่~\ref{tab:rag-conditions}

\begin{table}[ht]
  \centering
  \caption{เงื่อนไขการทดลองความจำเป็นของ RAG}
  \label{tab:rag-conditions}
  \begin{tabular}{|L{0.25\textwidth}|c|L{0.52\textwidth}|}
    \hline
    \textbf{การทดลอง} & \textbf{RAG} & \textbf{คำอธิบาย} \\ \hline
    1. Direct-Log      & ไม่ใช้ & ส่ง Representation Log ให้ LLM โดยตรง \\ \hline
    2. Direct-Summary  & ไม่ใช้ & ให้ LLM แปลง Log เป็นคำอธิบายพฤติกรรมภาษาคนก่อน แล้วจึงให้ระบุเทคนิค \\ \hline
    3.1 RAG-Log ($k=10$) & ใช้ & ใช้ Representation Log เป็น Query ค้นหา RAG ได้ Top-10 ให้ LLM \\ \hline
    3.2 RAG-Log ($k=5$)  & ใช้ & ใช้ Representation Log เป็น Query ค้นหา RAG ได้ Top-5 ให้ LLM \\ \hline
    4. RAG-Summary     & ใช้ & ใช้คำอธิบายพฤติกรรมภาษาคน (จากการทดลองที่ 2) เป็น Query ค้นหา RAG ได้ Top-$k$ ให้ LLM \\ \hline
  \end{tabular}
\end{table}

\subsection{ผลการทดลอง}
ผลการทดลองทั้ง 5 เงื่อนไข ได้แก่ อันดับของเฉลยในผลลัพธ์ของ RAG และคำตอบของ LLM แสดงในตารางที่~\ref{tab:rag-results}

\begin{table}[ht]
  \centering
  \caption{ผลการทดลองความจำเป็นของ RAG (คำตอบที่ถูกต้องคือ T1087.001)}
  \label{tab:rag-results}
  \small
  \begin{tabular}{|l|c|l|c|}
    \hline
    \textbf{การทดลอง} & \textbf{อันดับเฉลยใน RAG} & \textbf{คำตอบของ LLM} & \textbf{ผล} \\ \hline
    1. Direct-Log        & --                & T1083 File and Directory Discovery & ผิด \\ \hline
    2. Direct-Summary    & --                & T1078 Valid Accounts                & ผิด \\ \hline
    3.1 RAG-Log ($k=10$) & อันดับ 6          & T1087.001                           & ถูก \\ \hline
    3.2 RAG-Log ($k=5$)  & ไม่อยู่ใน Top-5   & T1033 System Owner/User Discovery   & ผิด \\ \hline
    4. RAG-Summary       & อันดับ 1          & T1087.001                           & ถูก \\ \hline
  \end{tabular}
\end{table}

\subsection{การวิเคราะห์ผล}
จากตารางที่~\ref{tab:rag-results} พบว่า
\begin{enumerate}
  \item เมื่อไม่ใช้ RAG (การทดลองที่ 1 และ 2) LLM ขนาดเล็กตอบผิดทั้งสองกรณี
        แม้ Input จะผ่านการคัดเลือก Event มาแล้ว แสดงว่าโมเดลขนาด 9B ยังขาดความรู้เฉพาะด้าน ATT\&CK
        ที่จะแยกแยะเทคนิคที่ใกล้เคียงกันได้
  \item LLM ตอบถูกเฉพาะเมื่อมี T1087.001 ซึ่งเป็นเฉลยอยู่ใน Top-$k$ ที่ได้จาก RAG
        ดังนั้นคุณภาพของ Retrieval เป็นปัจจัยสำคัญต่อความแม่นยำ
  \item การใช้คำอธิบายพฤติกรรมภาษาคนเป็น Query (RAG-Summary) ทำให้เฉลยขึ้นมาอยู่อันดับ 1
        ดีกว่าการใช้ JSON โดยตรง (อันดับ 6) เนื่องจากคำอธิบายภาษาคนมีความหมายใกล้เคียงกับคำอธิบายเทคนิคในฐานความรู้มากกว่า
  \item ค่า $k$ มีผลโดยตรง หาก $k$ น้อยเกินไปเฉลยอาจหลุดจาก Top-$k$ (การทดลอง 3.2)
\end{enumerate}

\subsection{ข้อจำกัดของการทดลอง}
การทดลองนี้ใช้เพียง 1 เทคนิคและ 1 กรณี และคำสั่ง \texttt{net users} เป็นคำสั่งที่มีลักษณะชัดเจน
จึงยังไม่สามารถสรุปผลโดยทั่วไปได้ จำเป็นต้องทดสอบกับชุดข้อมูลที่ครอบคลุมหลายเทคนิคเพิ่มเติม
และกำหนดเกณฑ์การประเมินไว้ล่วงหน้าก่อนดูผล

\section{แผนการทดลองและการประเมินผลในขั้นต่อไป}

\subsection{ชุดข้อมูล}
\begin{itemize}
  \item \textbf{Development Set} -- Splunk Attack Data กลุ่มที่มีร่องรอยการจำลองการโจมตีชัดเจน
        ใช้สำหรับออกแบบและปรับพารามิเตอร์ของ Pipeline
  \item \textbf{Test Set} -- OTRF Security-Datasets (ใช้ Adversary View เป็นข้อมูลอ้างอิง)
        และชุดข้อมูล Splunk ที่แยกไว้ไม่ใช้ระหว่างพัฒนา (Holdout)
\end{itemize}

\subsection{ตัวแปรในการทดลอง}
\begin{itemize}
  \item \textbf{Baseline} -- ส่ง Raw Logs (ภายใต้ขนาด Context เดียวกัน) ให้ LLM
  \item \textbf{Event Selection} -- เปรียบเทียบ 5 วิธี (Alert Only, Single Event, Time Window, Multiple Related, Top-$k$)
  \item \textbf{Representation} -- เปรียบเทียบ JSON, XML, TOON, Key=Value และ Timeline
  \item \textbf{RAG} -- ไม่ใช้ / ใช้ Query แบบ Log / ใช้ Query แบบ Summary และค่า $k$ ต่าง ๆ
\end{itemize}
โดยกำหนดให้ใช้ LLM ตัวเดียว (Qwen3.5:9B) และการตั้งค่าเดิมตลอดการทดลอง

\subsection{ตัวชี้วัด}
\begin{itemize}
  \item \textbf{Accuracy} -- ร้อยละของกรณีที่ระบุเทคนิคตรงกับ Ground Truth
        (แยกพิจารณาระดับ Sub-technique และระดับ Technique)
  \item \textbf{Precision / Recall / Macro F1} -- ความแม่นและความครอบคลุมรายเทคนิค
  \item \textbf{Retrieval Rank} -- อันดับของเฉลยในผลลัพธ์ของ RAG และการอยู่ใน Top-$k$
  \item \textbf{Latency} -- เวลาตั้งแต่รับ Logs จนได้คำตอบ (End-to-end)
  \item \textbf{Token Cost} -- จำนวน Token ที่ส่งให้ LLM ก่อนและหลังผ่าน Pipeline
\end{itemize}
